% Definitions section for "The Unchosen Adjacency" (revised: surplus and deficit)
% Assumes amsmath, amssymb, amsthm with definition/remark/proposition/proof environments.

\section{Definitions}\label{sec:defs}

\subsection{Relations, signed changes, and actions}

\begin{definition}[Relations and signed changes]
Let $R_u$ be the set of platform-mediated relations involving a user $u$ whose state may be changed by the platform or by an action exposed through its interface: who may view, comment, reply, tag, or message; what is recommended to whom; which commercial content appears beside the user's material; whether a restriction carries a stated reason or an end; and how the user's activity is attributed. A \emph{signed change} is a pair $(r,\varepsilon)$ with $r\in R_u$ and $\varepsilon\in\{+,-\}$, meaning that $r$ becomes present ($+$) or absent ($-$). Let $\Delta_u=R_u\times\{+,-\}$, and let $\tau:R_u\to\mathcal{T}$ assign each relation a type.
\end{definition}

\begin{definition}[Interface and effect set]
Let $\mathcal{A}$ be the set of actions the platform exposes to $u$, and $\mathcal{A}^{*}\supseteq\mathcal{A}$ the set of actions that are technically feasible for the platform to expose. What $u$ can select is represented by $\mathcal{A}$, and what the platform can vary by $R_u$. Each action $a\in\mathcal{A}^{*}$ has an \emph{effect set} $E(a)\subseteq\Delta_u$, the signed changes the platform actually makes.
\end{definition}

\subsection{Surplus and deficit}

\begin{definition}[Surplus, deficit, and mismatch]
For an action $a\in\mathcal{A}^{*}$ and a set $S\subseteq\Delta_u$ of changes the user intends, define
\[
U(a,S)=E(a)\setminus S,\qquad
D(a,S)=S\setminus E(a),\qquad
M(a,S)=U(a,S)\cup D(a,S).
\]
The \emph{surplus} $U$ is what the action does beyond the intention, and the \emph{deficit} $D$ is what the intention required and the action does not do. For $a\in\mathcal{A}$ with intended set $S_a$ write $U_a$, $D_a$, and $M_a$. Write $U^{+}$ and $U^{-}$ for the elements of $U$ with sign $+$ and $-$.
\end{definition}

No assumption is made that $S_a\subseteq E(a)$. An interface can fail an intention in two independent ways: by overreach ($U_a\ne\varnothing$) and by underrealization ($D_a\neq\varnothing$).

\begin{definition}[Forms of mismatch]
An action $a$ exhibits \emph{unchosen bundling} if $U_a\neq\varnothing$. It exhibits \emph{compelled association} if $U_a^{+}\neq\varnothing$, and \emph{compelled dissociation} if $U_a^{-}\neq\varnothing$. It exhibits a \emph{shortfall} if $D_a\neq\varnothing$.
\end{definition}

Publishing a post that is displayed with an advertisement adds the relation ``$u$ is adjacent to this advertisement'' and is compelled association. Blocking a person to revoke commenting removes that person's view of $u$ and is compelled dissociation. A restriction that bars commenting but carries no stated reason or end, when the user intended both, is a shortfall.

\begin{remark}
No relation between surplus effects and intended changes is modeled beyond set difference. That a single placement may attach at once to authorship, publication, identity, and continued participation affects no verdict below, and so no incidence structure is defined.
\end{remark}

\subsection{Attribution, value, and importance}

Mismatch alone is not condemnatory, since every software action has incidental effects and imperfect realizations. Five functions select the mismatch that matters: attribution and association by observers, the user's value for a retained relation, the user's aversion to an added relation, and the user's importance for an intended change. Attribution is descriptive and depends on context. The parameters are normative and depend on the kind of claim at stake.

\begin{definition}[Attribution and association]
For an observer $o$, a signed change $b$, and interface and interpretive conditions $c$, let
\[
J_o(b,u\mid c)\in[0,1]
\]
be the degree to which $o$ attributes $b$ to the choice, approval, or normative stance of $u$, and let
\[
A_o(b,u\mid c)\in[0,1]
\]
be the degree to which $o$ links $b$ with $u$ or with $u$'s content, without necessarily attributing choice or approval to $u$. Each action $a$ has a \emph{standard presentation context} $c_a$, the way its result is ordinarily displayed, including what shares its frame, what is adjacent to it in space, and what follows it in the viewing sequence.
\end{definition}

\begin{remark}
Association is distinct from attribution. A viewer may know that a user did not choose an advertisement and still experience it as linked to the user's content, because it shares a frame with it or follows it immediately. $J_o$ records the first judgment and $A_o$ the second. Whether associative linkage without attributed agency should count as consequential is a normative question, addressed by the thresholds below.
\end{remark}

\begin{definition}[Claim types and parameters]
Let $\kappa(b)$ classify the kind of attribution at issue for a signed change $b$, for example commercial endorsement, interpersonal rejection, institutional approval, or mere causal contribution. For each type $\kappa$ fix
\begin{itemize}
  \item a class $\mathcal{O}_\kappa$ of \emph{qualified observers}, whose readings are counted for claims of that type, with a probability measure $\mu_\kappa$ on it;
  \item an attribution threshold $\theta_\kappa\in(0,1]$ and an association threshold $\theta^{A}_\kappa\in(0,1]$; and
  \item a proportion $\rho_\kappa\in(0,1]$.
\end{itemize}
For each relation type $t\in\mathcal{T}$ fix de minimis levels $\lambda_t>0$ for losses and shortfalls and $\lambda^{+}_t>0$ for aversion.
\end{definition}

\begin{definition}[Value, aversion, and importance]
For a removed relation $b=(r,-)$, let $V_u(r)\ge0$ be the value $u$ places on retaining $r$. For an added relation $b=(r,+)$ and a context $c$, let $W_u(r\mid c)\ge0$ be the cost to $u$ of acquiring or sustaining $r$ under $c$. For an intended change $s\in S$, let $I_u(s)\ge0$ be the importance $u$ places on $s$ being made.
\end{definition}

\begin{remark}
The parameters $\theta_\kappa$, $\theta^{A}_\kappa$, $\rho_\kappa$, $\lambda_t$, $\lambda^{+}_t$ and the choice of qualified observers are openly normative. They are not psychological constants, and no single value is defended for all types. Reasonableness of a reading enters through the choice of $\mathcal{O}_\kappa$ and the thresholds, while $J_o$ records only what observers do attribute. The essay does not estimate the parameters. Declaring them makes clear why one eccentric reading or a trivial inconvenience does not condemn an interface.
\end{remark}

\begin{definition}[Consequential mismatch]
Let $a\in\mathcal{A}^{*}$ and $S\subseteq\Delta_u$, and let $c_a$ be the standard presentation context of $a$. An element of $M(a,S)$ is \emph{consequential} if it satisfies one of the following.
\begin{enumerate}
  \item (\emph{misattribution}) It is some $b\in U(a,S)$ with $\kappa=\kappa(b)$ and
  \[
  \mu_\kappa\bigl(\{\,o\in\mathcal{O}_\kappa : J_o(b,u\mid c_a)\ge\theta_\kappa\,\}\bigr)\ \ge\ \rho_\kappa .
  \]
  \item (\emph{association}) It is some $b\in U(a,S)$ with $\kappa=\kappa(b)$ and
  \[
  \mu_\kappa\bigl(\{\,o\in\mathcal{O}_\kappa : A_o(b,u\mid c_a)\ge\theta^{A}_\kappa\,\}\bigr)\ \ge\ \rho_\kappa .
  \]
  \item (\emph{loss}) It is some $b=(r,-)\in U(a,S)$ with $V_u(r)\ge\lambda_{\tau(r)}$.
  \item (\emph{aversion}) It is some $b=(r,+)\in U(a,S)$ with $W_u(r\mid c_a)\ge\lambda^{+}_{\tau(r)}$.
  \item (\emph{shortfall}) It is some $s=(r,\varepsilon)\in D(a,S)$ with $I_u(s)\ge\lambda_{\tau(r)}$.
\end{enumerate}
\end{definition}

\begin{remark}
Aversion complements attribution and loss and does not replace them. An added relation can be consequential because observers read it as the user's act, or because the user objects to sustaining it, for instance by contributing content, attention, or legitimacy to a system the user rejects, even when no observer reads it as endorsement. The threshold $\lambda^{+}_t$ prevents every disliked side effect from becoming consequential. Establishing consequence does not establish a defect, since feasibility and unavailability remain to be shown. The definitions are stated for a single user; platform-level evaluation would require a distribution over user types, which the essay leaves for later work.
\end{remark}

\subsection{Realizability and defective actions}

\begin{definition}[Clean realization]
An intended set $S\subseteq\Delta_u$ is \emph{cleanly realizable at} an interface $\mathcal{B}$ if there is $a'\in\mathcal{B}$ such that $M(a',S)$ contains no consequential element. Exact realization $E(a')=S$ is the special case with empty mismatch.
\end{definition}

\begin{remark}
Exact realization is too strong a requirement, since a benign log entry would defeat it. Cleanliness is relative to consequence, and consequence is defined through $J_o$, $V_u$, and $I_u$ without reference to realizability, so the definition is not circular. An intended set for which no action in $\mathcal{B}$ makes the required changes is simply not cleanly realizable there.
\end{remark}

\begin{definition}[Withheld intention and defective action]
An intended set $S$ is \emph{withheld} at $\mathcal{A}$ if it is cleanly realizable at $\mathcal{A}^{*}$ but not at $\mathcal{A}$. An action $a\in\mathcal{A}$ is \emph{defective} if $S_a$ is withheld at $\mathcal{A}$. A defective action is a \emph{defective bundle} if a consequential element of $M_a$ lies in $U_a$, and a \emph{defective shortfall} if one lies in $D_a$.
\end{definition}

Because $a$ itself is available and does not realize $S_a$ cleanly, a defective action has at least one consequential element of $M_a$, so it is a defective bundle, a defective shortfall, or both. The withheld clause is the counterfactual criterion. It condemns an interface only when the platform could have offered an action that makes the intended changes without consequential mismatch and did not. Unavoidable side effects and unavoidable limits are excluded.

\subsection{Exit}

\begin{proposition}[Cost of the intended change]
If $S$ is withheld at $\mathcal{A}$, then every $a'\in\mathcal{A}$ has a consequential element in $M(a',S)$.
\end{proposition}

\begin{proof}
Otherwise some $a'\in\mathcal{A}$ would cleanly realize $S$ at $\mathcal{A}$, contradicting that $S$ is withheld.
\end{proof}

The user can therefore obtain $S$ only by accepting some consequential mismatch, surplus or deficit, which need not be that of the original action. The only way to avoid consequential mismatch altogether is to forgo $S$, by refraining or leaving the platform. This is a condition on the interface, not a claim of legal compulsion. The availability of exit is not evidence that a clean realization exists.
